\section{How the framework composes across mathematical domains}
\label{sec:applications}

The simulation component and the finite-algebra component have different
jobs. Simulation preserves a proved observation during a finite
construction. Finite observable modules restrict the possible rational
values of an observation. Reconstruction turns the resulting finite
possibilities into actual objects when a finite fiber procedure is
available. The following examples show how to use the three steps
without adding an implicit existence principle.

\subsection{Modular computation and exact polynomial evaluation}

Fix $q\ge1$ and let $\rho_q(a)$ be the remainder of $a$ on division
by $q$. Define the source step $f(a)=a^2$ and the target step
$g(r)=\rho_q(r^2)$. The local fusion law is
\[
                  \rho_q(a^2)=\rho_q(\rho_q(a)^2).
\]
To prove it arithmetically, write $a=qA+r$, $0\le r<q$.
Expanding gives $a^2=q(qA^2+2Ar)+r^2$. Division of $r^2$ by
$q$ therefore gives the same final remainder as division of $a^2$;
uniqueness follows by subtracting two division equations with remainders
between zero and $q-1$. These are finite integer equations and
inequalities, with a quantifier-free induction proof for division.
Theorem~\ref{thm:fusion} now transports the entire repeated-squaring
construction through $\rho_q$.

The statement is uniform in the numerical modulus. To instantiate the
unary theorem literally, use state $(q,a)$ and a representation map
$(q,a)\mapsto(q,\rho_q(a))$. The guard $q\ge1$ belongs in the
relational invariant; a total default for $q=0$ is never used to
justify modular mathematics. This pattern applies equally to computations
in a presented number ring after its finite arithmetic and reduction
laws have been supplied.

For a rational polynomial $P$, evaluation at a fixed rational tuple
$p$ preserves addition and multiplication. The multiplication assertion
has a direct finite proof: distribute the two coefficient sums, multiply
the monomials, and use $p^{\alpha+\beta}=p^\alpha p^\beta$.
There are only the finitely many coefficient pairs in the input lists.
The local proof transports repeated squaring and more general finite
expression evaluation into rational arithmetic. It can also transport
coefficient identities into equations at a common zero, the precise
bridge used by the certificate theorem.

\subsection{Truncated series and coherent exact computation}

Let $\pi_N$ discard all coefficients above degree $N$. For finite
polynomials or coefficient lists, convolution gives
\[
      \pi_N(AB)=\pi_N\bigl(\pi_N(A)\pi_N(B)\bigr).
\]
Indeed the coefficient of degree $k\le N$ is
$\sum_{i=0}^{k}a_i b_{k-i}$. Every contributing index is at most
$N$, so truncating either input leaves it unchanged. For $k>N$
both truncated outputs have zero coefficient. Coefficient equality
is a bounded list condition and hence an eligible arithmetic invariant.

Iteration fusion applies to repeated polynomial operations whose local
truncation law has this form. It yields the exact requested jet without
requiring a completed infinite series as an object of \ipra. This
connects to the repository's filtered-lifting and formal-series work
\cite{Coherent}. Differentiation uses a different demand: coefficients
through $N+1$ determine the derivative through $N$. A correct
simulation must include that shift rather than claiming that arbitrary
operations preserve the same truncation level.

Exact coherence does not imply an analytic tail bound. A convergent
analytic interpretation must supply a PR function $T(N,k)$ or another
explicit estimate making the omitted tail small enough at requested
precision $k$. The present transport theorem then handles the finite
coefficient calculation; the tail theorem remains its own mathematical
input. This is a common interface for the Gamma, zeta, calculus,
Lagrange inversion and regular local-branch developments, rather than
a replacement for their analytic estimates.

\subsection{Redundant algebraic presentations and observational equality}

The point-enumeration checker intentionally does not require a basis or
commuting matrices. This is useful, not just a technical weakening.
Consider
\[
 F_1=x^2-x,\qquad F_2=2y-x,\qquad (b_1,b_2,b_3)=(1,x,x).
\]
Use
\[
 M_x=\begin{pmatrix}0&0&0\\1&1&1\\0&0&0\end{pmatrix},
 \qquad
 M_y=\begin{pmatrix}0&0&0\\0&0&0\\1/2&1/2&1/2\end{pmatrix}.
\]
They do not commute: the nonzero row of $M_xM_y$ is its second,
whereas that of $M_yM_x$ is its third. Yet the action identities
are valid. For the $x$ action the residuals are $0,F_1,F_1$.
For the $y$ action they are
\[
 \tfrac12F_2,\qquad
 \tfrac12F_1+\tfrac x2F_2,\qquad
 \tfrac12F_1+\tfrac x2F_2.
\]
The unit generator is $b_1=1$. Hence this is an accepted certificate.
It gives $\Delta=(1,2)$ and $R=(1,1)$, and the grid filter returns
exactly $(0,0)$ and $(1,1/2)$.

The apparent matrix discrepancy acts in a redundant direction. At any
common zero, vectors of coefficients $c,d$ representing the same
value are related by
\[
                  \sum_i c_i b_i(p)=\sum_i d_i b_i(p).
\]
That relation, with the common-zero guard and the fixed point $p$
stored as parameters, is the right observation to preserve. Equality
of coefficient vectors would be unnecessarily strong. A relational
simulation therefore allows a convenient redundant presentation without
first constructing a canonical basis. This matters in weak arithmetic,
where basis selection, rank certificates and coefficient growth are
additional obligations which the desired conclusion may not need.

\subsection{A finite open part of a positive-dimensional system}

Equations plus guards occur in rational maps, affine charts, and
elimination with nonzero denominators. Here is a complete concrete
example of the localization principle.
The equations
\[
                xy=0,\qquad x(x-1)=0
\]
have infinitely many rational solutions: the entire line $x=0$ is
included. On the open part $x\ne0$, the only solution is $(1,0)$.
Introduce $t$ and use
\[
              F_1=xy,\quad F_2=x(x-1),\quad F_3=tx-1.
\]
The one-element spanning family $(1)$ has action matrices $M_x=(1)$,
$M_y=(0)$, $M_t=(1)$. The complete certificate is
\begin{align*}
 x-1&=tF_2-(x-1)F_3,\\
 y&=tF_1-yF_3,\\
 t-1&=-t^2F_2+\bigl(1+t(x-1)\bigr)F_3.
\end{align*}
Each identity follows by expanding both sides; together with the unit
identity they are all the checker needs. The coordinate radii are
$(1,0,1)$, so nine candidate triples suffice. Filtering gives precisely
$(1,0,1)$; projection gives the desired open-part list.

This example illustrates a concrete certificate producer in a useful
special case. It does not infer that localization always makes a system
finite. The guard is explicit, its inverse coordinate is unique, and
the extended system has its own verified action data. Several guards
$G_i\ne0$ can be combined as $G=\prod_iG_i\ne0$ over $\Q$,
or treated with separate inverse coordinates.

\subsection{Why finite rational sets on curves are a different problem}

A multiplication certificate is stronger than the assertion that the
rational zero set is finite. The certificate identities imply that the
entire quotient algebra $\Q[X_1,\ldots,X_r]/(F_1,\ldots,F_s)$
is spanned by finitely many classes: start with the unit combination,
multiply by variables using the action tables, and extend by linearity.
This argument applies to all polynomials, without testing rational points.

For a single monic plane equation $F(X,Y)$ of positive degree in $Y$,
the map $\Q[X]\to\Q[X,Y]/(F)$ is injective. A nonzero multiple
$HF$ has $Y$-degree $\deg_Y H+\deg_Y F>0$ and so cannot be a
nonzero polynomial in $X$ alone. Consequently the powers of $X$
remain linearly independent in the quotient. No finite spanning
certificate for both coordinates can exist for that full ideal.
This remains true if the curve has only finitely many rational points.

Thus the present theorem does not furnish a direct decision procedure
for rational points on every curve. It applies to finite algebraic
sections, suitably finite open charts, or other situations for which
the extra finite action data are actually available. The norm constructor
explains exactly how an additional section equation supplies those data.
The missing global arithmetic work is to show that the points of
interest are covered by finitely many such explicit sections, with
uniformly available coefficients or a suitable bound on them.

\subsection{A precise task decomposition for the difficult reports}

\begin{longtable}{P{.20\linewidth}P{.38\linewidth}P{.32\linewidth}}
\toprule
Area & Reusable interface supplied here & Further theorem still required\\
\midrule\endhead
Siegel-style sections & A nonzero norm produces all rational points
in a supplied monic-plane section, and a finite algebra certificate. &
A global cover or bounded coefficients for its sections; the quantitative
approximation theorem is not supplied.\\[4pt]
Faltings-style finiteness & Complete lists for individually certified
finite algebraic pieces; finite union and chart transport. & A method
covering the rational points of a positive-dimensional curve by such
explicit pieces.\\[4pt]
Mordell--Weil descent & Proof transport for finite descent computations
and exact enumeration of separately certified finite loci. & Actual
covering representatives and a proof that they cover the required group
quotient.\\[4pt]
Formal algebra and local expansions & Checked simulation through
truncation and evaluation, once their local object laws are supplied. &
Application-specific local laws, singular branch analysis, and any
analytic tail bounds.\\[4pt]
Analytic number theory & Composable finite exact arithmetic and
explicit certificate checking. & The uniform analytic estimates named
in the prime-distribution report.\\
\bottomrule
\end{longtable}

This division is productive because it names executable and provable
subtasks. A future global-cover theorem can call the enumeration module
directly. A future formal coefficient library can reuse the checked
simulation builder. Until those inputs exist, the completed sublemmas
remain useful without being misreported as proofs of the original
deep theorems.
